\subsection{Double Integrals in Polar Coordinates}


\content{

}{% presentation
\vspace{3in}
}{% nopresentation
}{% comments
Review polar coordinates, go over how the extra r shows up in the integration:
\begin{enumerate}
\item draw a basic region
\item draw sectors based on $\Delta \theta$ and $\Delta r$
\item show that the area of each polar rectangle is given by 
$$\Delta A_i = \frac{1}{2}(r_i^2-r_{i-1}^2)\Delta \theta$$
by using the difference of the areas of two sectors of a circle
\item use the fact that $r_i^* = \frac{1}{2}(r_i+r_{i-1})$ is the center point
of each polar rectangle to show that the area is 
$$ \Delta A_i = r_i^* \Delta r \Delta \theta$$
\end{enumerate}
}

\content{
\begin{theorem} If $f$ is a continuous function on a polar rectangle $R$ given
by $a\leq \theta\leq b$ and $\alpha\leq r\leq \beta$, then 
$$ \iint_Rf(x,y)\,dA = \int_\alpha^\beta\int_a^b rf(r\cos\theta,
r\sin\theta)\,d\theta dr. $$
\end{theorem}
}{% presentation
}{% nopresentation
}{% comments
}

\content{
\begin{example} Evaluate the integral 
$$ \iint_R (3x+4y^2)\,dA $$
where $R$ is the region in the upper-half plane bounded by the circles
$x^2+y^2=1$ and $x^2+y^2=4$.
\end{example}
}{% presentation
\vspace{4in}
}{% nopresentation
}{% comments
}

\content{
\begin{example} Find the volume of the solid bounded by the plane $z=0$ and the
surface $z=1-x^2-y^2$.
\end{example}
}{% presentation
\vspace{3in}
}{% nopresentation
}{% comments
}

\content{
\begin{example} Use a double integral to find the area of the region enclosed by
one loop of the function $r=\cos 2\theta$.
\end{example}
}{% presentation
\vspace{4in}
}{% nopresentation
}{% comments
}

\content{
\begin{example} FInd the volume of the solid that lies in the cylinder
$x^2+y^2=2x$ and between the $xy$-plane and the surface $z=x^2+y^2$.
\end{example}
}{% presentation
\vspace{4in}
}{% nopresentation
}{% comments
Complete the square in the cylinder equation to find it's center and radius
then set $x=r\cos\theta$ and substitute this into the cylinder equation (and
$r^2=x^2+y^2$)

Mention that we will be skipping 15.4, which contains many applications of
double integrals. 
}
